\documentclass[11pt]{article}
\usepackage[margin=1in]{geometry}
\usepackage{amsmath,amssymb,amsthm,mathtools}
\usepackage{graphicx,xcolor,microtype,hyperref}
\hypersetup{colorlinks=true,linkcolor=blue!55!black,citecolor=blue!55!black,urlcolor=blue!55!black}
\newtheorem{theorem}{Theorem}
\newtheorem{lemma}{Lemma}
\newcommand{\R}{\mathbb R}
\newcommand{\Z}{\mathbb Z}
\newcommand{\G}{\Gamma}
\newcommand{\AEn}{\mathsf{AE}}
\newcommand{\W}{\mathsf W_1}
\newcommand{\norm}[1]{\left\lVert#1\right\rVert}
\title{\vspace{-2em}Uniform Transport F{\o}lner Sets\\for Virtually Abelian Groups}
\author{Partial solution packet for arXiv:2509.10686}
\date{17 August 2026}

\begin{document}
\maketitle
\begin{center}
\fcolorbox{orange!75!black}{orange!4}{\parbox{0.91\textwidth}{
\textbf{Status: candidate substantial partial theorem, likely valid.}
Problem~10.8 has an affirmative answer for every finitely generated virtually
abelian group satisfying the source's no-real-characters hypothesis.  Explicit
lattice-box sets have transport defect $O(1/N)$.  The general amenable-group
problem remains open.
}}
\end{center}

\section{Source problem and new result}

For a finite nonempty $F\subseteq\G$, write
\[
                       u_F=\frac1{|F|}\sum_{x\in F}\delta_x .
\]
On a finitely generated group with a left-invariant word metric, the
Wasserstein distance between finitely supported probabilities equals the
Arens--Eells norm of their difference:
\[
                         \W(\mu,\nu)=\norm{\mu-\nu}_{\AEn}.
\]
Rosendal proves that if $\G$ is amenable and has no nonzero homomorphism to
$\R$, then, for every finite $E\subseteq\G$ and $\varepsilon>0$, there is a
finitely supported probability $\beta$ with
\[
             \max_{g,f\in E}\W(\beta g,\beta f)<\varepsilon
             \tag{1}
\]
\cite[Theorem~10.4]{R}.  Rational approximation makes $\beta$ a uniform
average of a finite \emph{multiset}.  Problem~10.8 asks whether repetitions
can always be removed, so that $\beta=u_F$ for a genuine finite set $F$.

\begin{figure}[ht]
\centering
\includegraphics[width=0.94\textwidth]{figures/open_problem_crop.png}
\caption{The motivation and Problem~10.8 on source page~32.  The page was
reconstructed from the run's ingested official arXiv source.}
\end{figure}

We prove this for the full virtually abelian class.

\subsection*{Idea of the proof}
Use every coset of a finite-index lattice subgroup at once.  Right
translation shifts the lattice box in each coset by a bounded displacement.
The sum of those displacements is a homomorphism from $\G$ to the lattice, so
the no-real-characters hypothesis forces its value to vanish.  After moving
the small boundary defects into one coset layer, the resulting signed lattice
kernel has zero mass and zero first moment.  Algebraically it is therefore a
sum of second differences.  A second difference of a normalized box has
Arens--Eells norm $O(1/N)$, which is the required transportation estimate.
The argument does not provide such sets for general amenable groups.

\begin{theorem}\label{thm:main}
Let $\G$ be a finitely generated virtually abelian group with no nonzero
homomorphism $\G\to\R$, and equip it with any left-invariant word metric.
There are finite nonempty sets $F_N\subseteq\G$ such that, for every
$g\in\G$,
\[
                            \W(u_{F_N}g,u_{F_N})=O_g(N^{-1}).
                            \tag{2}\label{eq:rate}
\]
Consequently, for every finite $E\subseteq\G$ and every $\varepsilon>0$, some
$F_N$ satisfies
\[
                    \max_{g,f\in E}\W(u_{F_N}g,u_{F_N}f)<\varepsilon.
\]
\end{theorem}

The infinite dihedral example in the source is the rank-one case.  The proof
below identifies its opposite-drift cancellation as a general transfer
identity.

\section{Coset displacement has zero total drift}

If $\G$ is finite, take $F_N=\G$.  Otherwise choose a normal finite-index
torsion-free abelian subgroup $A\cong\Z^d$.  Such an $A$ exists after taking
the normal core of a torsion-free finite-index abelian subgroup.  Fix a set
$T$ of representatives, so every $x\in\G$ is written uniquely as
\[
                               x=a t,\qquad a\in A,\quad t\in T.
\]
Write $A$ additively and put $m=|T|$.

For $g\in\G$ and $t\in T$, there are unique
$c_t(g)\in A$ and $\sigma_g(t)\in T$ such that
\begin{equation}\label{eq:cocycle}
                         t g=c_t(g)\,\sigma_g(t).
\end{equation}
Here $\sigma_g$ is the permutation of the coset space induced by right
multiplication.

\begin{lemma}[Transfer cancellation]\label{lem:transfer}
For every $g\in\G$,
\[
                              \sum_{t\in T}c_t(g)=0\quad\text{in }A.
\]
\end{lemma}

\begin{proof}
From \eqref{eq:cocycle},
\[
 c_t(gh)=c_t(g)+c_{\sigma_g(t)}(h).
\]
Since $\sigma_g$ permutes $T$, the map
\[
                         S(g)=\sum_{t\in T}c_t(g)
\]
is a group homomorphism $\G\to A$.  If $S(g)\neq0$ for some $g$, a real
linear functional on $A\otimes_{\Z}\R$ nonzero at $S(g)$ would make a
nonzero homomorphism $\G\to\R$.  This contradicts the hypothesis, so $S=0$.
\end{proof}

\section{Box estimates and layer collapse}

Fix a basis $e_1,\ldots,e_d$ of $A$ and put
\[
 B_N=\left\{\sum_{i=1}^d n_ie_i:0\leq n_i<N\right\},
 \qquad b_N=\frac1{|B_N|}\sum_{a\in B_N}\delta_a,
 \qquad F_N=B_NT.
\]
For $v\in A$, let $\tau_v$ denote translation by $v$.  A coordinatewise
boundary count gives
\begin{equation}\label{eq:tvbox}
 \norm{\tau_vb_N-b_N}_1\leq \frac{2\sum_i|v_i|}{N}
 \qquad (v=\textstyle\sum_i v_ie_i),
\end{equation}
with the harmless convention that the right side may exceed $2$.

For $s\in T$, write $\iota_s\eta$ for a finitely supported signed measure
on the layer $As$ obtained from a signed measure $\eta$ on $A$.  Fix
$t_0\in T$.  If $\eta(A)=0$, then
\begin{equation}\label{eq:collapse}
 \norm{\iota_s\eta-\iota_{t_0}\eta}_{\AEn}
 \leq d(s,t_0)\norm{\eta}_1.
\end{equation}
Indeed, the left side is represented by
$\sum_a\eta(a)(\delta_{as}-\delta_{at_0})$, and left invariance gives
$d(as,at_0)=d(s,t_0)$.

For fixed $g\in\G$, set
$d_s=c_{\sigma_g^{-1}(s)}(g)$.  Right translating $F_N$ and then relabelling
the coset layers gives the exact identity
\begin{equation}\label{eq:defect}
 u_{F_N}g-u_{F_N}
 =\frac1m\sum_{s\in T}
   \iota_s\bigl(\tau_{d_s}b_N-b_N\bigr).
\end{equation}
Each summand has mass zero.  Thus \eqref{eq:tvbox} and
\eqref{eq:collapse} show that replacing every layer $s$ in
\eqref{eq:defect} by the fixed layer $t_0$ changes the Arens--Eells norm by
$O_g(N^{-1})$.

\section{The zero-drift kernel is a second difference}

Define the finitely supported signed measure
\[
 \kappa_g=\frac1m\sum_{s\in T}(\delta_{d_s}-\delta_0)
 \quad\text{on }A.
\]
It has total mass zero.  Lemma~\ref{lem:transfer} also gives zero first
moment:
\[
 \sum_{a\in A}a\,\kappa_g(a)=\frac1m\sum_{s\in T}d_s=0.
\]

\begin{lemma}[Second-difference decomposition]\label{lem:square}
If a finitely supported $\kappa\in\R[A]$, $A\cong\Z^d$, has zero mass and
zero first moment, then there are finitely many
$\alpha_r\in\R$ and $v_r,w_r\in A$ such that
\[
 \kappa=\sum_r\alpha_r
 (\delta_{v_r}-\delta_0)*(\delta_{w_r}-\delta_0).
\]
\end{lemma}

\begin{proof}
Let $I$ be the augmentation ideal of the Laurent group algebra $\R[A]$.
The zero-mass condition says $\kappa\in I$.  Modulo $I^2$ one has
\[
 \delta_v-\delta_0\equiv
 \sum_{i=1}^d v_i(\delta_{e_i}-\delta_0)
 \qquad(v=\textstyle\sum_i v_ie_i),
\]
also for negative coordinates, by the Laurent identities.  Hence the class
of $\kappa$ in $I/I^2$ is exactly its first-moment vector.  That vector is
zero, so $\kappa\in I^2$.  Since $I$ is generated by the elements
$\delta_v-\delta_0$, membership in $I^2$ gives the asserted finite
decomposition.
\end{proof}

For fixed $v,w\in A$, let $\nu_N=\tau_wb_N-b_N$.  Pairing each coefficient
of $\nu_N$ with its translate by $v$ inside the layer $At_0$ gives
\begin{equation}\label{eq:second}
 \norm{\iota_{t_0}(\tau_v\nu_N-\nu_N)}_{\AEn}
 \leq d(vt_0,t_0)\norm{\nu_N}_1
 =O_w(N^{-1}),
\end{equation}
where the distance is a fixed constant depending only on $v,t_0$ and the
chosen word metric.  Applying Lemma~\ref{lem:square} to $\kappa_g$ and
summing \eqref{eq:second} yields
\begin{equation}\label{eq:collapsed}
 \norm{\iota_{t_0}(b_N*\kappa_g)}_{\AEn}=O_g(N^{-1}).
\end{equation}

\section{Proof of the theorem}

\begin{proof}[Proof of Theorem~\ref{thm:main}]
The layer-collapsed form of \eqref{eq:defect} is exactly
$\iota_{t_0}(b_N*\kappa_g)$.  The collapse error is $O_g(N^{-1})$ and
\eqref{eq:collapsed} has the same bound.  Therefore
$\W(u_{F_N}g,u_{F_N})=O_g(N^{-1})$.

For finite $E$, the triangle inequality gives
\[
 \W(u_{F_N}g,u_{F_N}f)
 \leq \W(u_{F_N}g,u_{F_N})+
       \W(u_{F_N},u_{F_N}f),
\]
which tends to zero uniformly over the finitely many pairs $g,f\in E$.
Choosing $N$ sufficiently large completes the proof.
\end{proof}

\section{Verification, scope, and novelty}

The proof was independently reconstructed from the source problem.  The
three points most vulnerable to hidden assumptions were checked explicitly:
the cocycle sum is a genuine homomorphism because $A$ is abelian; moving a
zero-mass signed boundary measure between coset layers costs its
$\ell_1$-mass times a fixed layer distance; and the augmentation-ideal
quotient $I/I^2\cong A\otimes\R$ identifies the remaining obstruction with
the first moment.

No computation is used as proof.  The crop script was rerun from the local
source rendering.  A bounded novelty search on 21 August 2026 covered the
run indexes, exact title and arXiv-id searches, exact phrases from
Problem~10.8, and primary searches combining uniform finite subsets with the
Wasserstein/Arens--Eells formulation.  No later resolution or statement of
this virtually abelian theorem was located.

The result is only partial: it settles the full virtually abelian class under
the source hypothesis, but the general finitely generated amenable case and
even a virtually nilpotent extension are not proved here.

\paragraph{Human review recommendation.}
Check identities \eqref{eq:defect} and \eqref{eq:collapse}, the identification
of the first moment with the class in $I/I^2$, and the second-difference
transport bound \eqref{eq:second}.

\begin{thebibliography}{9}
\bibitem{R}
C. Rosendal,
\emph{Amenability, Optimal Transport and Abstract Ergodic Theorems},
arXiv:2509.10686v1 (2025).
\end{thebibliography}

\end{document}
